\documentclass[12pt]{article}
\usepackage[utf8]{inputenc}
\usepackage[T1]{fontenc}
\usepackage{hyperref}
\title{Southern Tutchone Language Resources for K-12 Curriculum Development}
\author{Compiled by Claude Sonnet 5 (claude-sonnet-5) for Andr\'as Kornai}
\date{}
\begin{document}
\maketitle
\section{Introduction}
Southern Tutchone (D\"an K\'e; ISO 639-3 tce) is an Athabaskan language of the Yukon, Canada, with a small number of speakers. It has the best-documented school presence of the four languages in this batch.
\section{Primary Educational Resources}
Search summaries report that Elias Community School in Haines Junction teaches Southern Tutchone from K to 12, that a bicultural Southern Tutchone kindergarten programme began there in 2009, and that the Yukon Native Language Centre (YNLC) supports a ``Dakeyi -- Our Country'' programme for high school students. These come from the search summary, not a page I opened. Yukon University has offered a beginner course (2016 outline):
\url{https://www.yukonu.ca/sites/default/files/course-outlines/F_16_ATHA101_Bunbury_Instructing_Outline_Prepped_byAllison_.pdf}
\section{Linguistic and Reference Materials}
YNLC English-to-Southern Tutchone noun dictionary arranged by subject (compiled by Margaret Workman, 2002 draft; PDF with audio): \url{https://ynlc.ca/wp-content/uploads/st_dictionary_englishtost_bysubject.pdf}

Catalogue record: \url{https://wiki.ubc.ca/Documentation:RelLex/Southern_Tutchone_Noun_Dictionary}

YNLC language overview: \url{https://ynlc.ca/?p=4033}

Wikipedia: \url{https://en.wikipedia.org/wiki/Southern_Tutchone_language}
\section{Religious and Community Resources}
CBC report on a Southern Tutchone dictionary created by Yukon sisters: \url{https://www.cbc.ca/lite/story/1.6706168}
\section{Digital and Audio Resources}
In 2018 YNLC digitized at least 10,000 pages of material across eight Yukon languages and made it searchable online: \url{https://www.cbc.ca/lite/story/1.4761056}
\section{Curriculum Development Context}
Existing school instruction is as a subject. Dictionary and digitized materials give a base, but speakers are few and elderly, so immersion-oriented materials and audio recording are priorities.
\end{document}